% Part: lambda-calculus
% Chapter: syntax
% Section: terms

\documentclass[../../../include/open-logic-section]{subfiles}

\begin{document}

\olfileid{lam}{syn}{trm}
\olsection{ترمونه}

د لامبډا حساب ترمونه د متغيرونو له نامتناهي زېرمي $\Obj{v_0}$،
$\Obj{v_1}$، \dots، له~``$\lambd$'' نښې او قوسونو څخه په استقرايي
ډول جوړېږي۔ د متغيرونو د ښودلو دپاره $x$، $y$، $z$، \dots{} او
د ترمونو د ښودلو دپاره $M$، $N$، $P$، \dots{} کاروو۔

\begin{defn}[ترمونه] \ollabel{defn:term}
د لامبډا حساب د \emph{ترمونو} سټ په استقرايي ډول داسې تعريفېږي:
\begin{enumerate}
  \item \ollabel{defn:term-var} که $x$ متغير وي، نو $x$ ترم دے۔
  \item \ollabel{defn:term-abs} که $x$ متغير او $M$ ترم وي، نو
    $(\lambd[x][M])$ ترم دے۔
  \item \ollabel{defn:term-app} که $M$ او $N$ دواړه ترمونه وي، نو
    $(MN)$ ترم دے۔
\end{enumerate}
\end{defn}

که ترم $(\lambd[x][M])$ د \olref{defn:term-abs} له مخې جوړ شوے
وي، نو وايو چې دا د \emph{تجريد} پايله ده، او په $\lambd[x]$ کښې
$x$ ته \emph{!!{parameter}} وايو۔ که ترم $(MN)$ د
\olref{defn:term-app} له مخې جوړ شوے وي، نو دا د \emph{تطبيق}
پايله ده۔

پاس تعريف شوي ترمونه ټولې قوسې په بشپړه توګه لري۔ دا کار يو څه
دروند کېدے شي، لکه ترم
$(\lambd[x][((\lambd[x][x])(\lambd[x][(xx)]))])$ چې ښيي۔ د
قوسونو د پرېښودلو دپاره به نښيز دودونه وټاکو۔ خو رسمي تعريف دا
آسانوي چې د \olref{defn:term} له مخې د هر ترم د جوړېدو طريقه
وپېژنو۔ د بېلګې په توګه، د ترم
$(\lambd[x][((\lambd[x][x])(\lambd[x][(xx)]))])$ د جوړېدو وروستے
ګام بايد تجريد وي چې !!{parameter} ئې~$x$ دے۔ دا ترم له
$((\lambd[x][x])(\lambd[x][(xx)]))$ څخه د تجريد په وسيله جوړېږي؛
او دغه ترم د دوو ترمونو تطبيق دے۔ له دغو دوو څخه هر ترم بيا د
تجريد پايله ده، او همداسې نور۔

\begin{prob}
د $(\lambd[g][(\lambd[x][(g (x x))])
  (\lambd[x][(g (x x))])])$ د جوړېدو طريقه بيان کړئ۔
\end{prob}

\end{document}
